\problema{Kth Largest Element in an Array}{215}{src/04-heaps-topk/215-kth-largest-element.cpp}{M}

Nos dan un arreglo \texttt{nums} sin ordenar y un entero \texttt{k}. Hay que
devolver el \texttt{k}-ésimo elemento más grande, contando desde el más
grande (el primer más grande es el máximo, el segundo es el que le sigue, y
así sucesivamente). Si hay valores repetidos, cada copia cuenta por
separado: no es el \texttt{k}-ésimo valor \emph{distinto}.

\paragraph{La idea, con un ejemplo de la vida real.} Imagina que estás
organizando una competencia y solo te interesa saber quién quedó en el
\texttt{k}-ésimo lugar del podio, sin importarte el resto de la tabla
completa. No hace falta ordenar a todos los competidores de principio a fin:
basta con quedarte, en todo momento, con los \texttt{k} mejores vistos hasta
ahora, y descartar de inmediato a cualquiera que no logre meterse en ese
grupo selecto.

Esa es la idea de usar un \emph{min-heap} (montículo mínimo, una estructura
que siempre te da acceso rápido a su elemento más pequeño) de tamaño fijo
\texttt{k}: metemos cada número, y si el montículo crece más de \texttt{k}
elementos, sacamos el más chico de ese grupo, porque ya sabemos que no puede
ser el \texttt{k}-ésimo más grande de todo el arreglo.

\begin{center}
\ttfamily
\begin{tabular}{l}
nums = [3,2,1,5,6,4], k = 2\\
ordenado desc: 6, 5, 4, 3, 2, 1 -> 2do lugar = 5\\
\end{tabular}
\end{center}

\paragraph{Cómo lo resuelve el código, paso a paso.}
\begin{enumerate}
  \item Creamos un min-heap vacío.
  \item Por cada número del arreglo, lo insertamos en el heap.
  \item Si el heap tiene más de \texttt{k} elementos, sacamos el más
        pequeño: ya perdió toda oportunidad de ser el \texttt{k}-ésimo mayor.
  \item Al terminar de recorrer el arreglo, el heap contiene exactamente los
        \texttt{k} valores más grandes vistos, y como es un min-heap, su
        cima (el menor de esos \texttt{k}) es justo la respuesta.
\end{enumerate}

\lstinputlisting{src/04-heaps-topk/215-kth-largest-element.cpp}

\complejidad{$O(n \log k)$}{$O(k)$}

\paragraph{¿Qué significa esa complejidad?} Recorremos los $n$ elementos del
arreglo una vez, y cada operación sobre el heap (insertar o sacar) cuesta
$\log k$ porque el heap nunca guarda más de \texttt{k} elementos a la vez.
Como normalmente \texttt{k} es mucho más pequeño que $n$, esto es bastante
más rápido que ordenar todo el arreglo ($O(n \log n)$), sobre todo cuando
\texttt{k} es chico.

\paragraph{Consejos para la entrevista (Amazon).} Este problema suele usarse
para ver si sabes elegir la estructura correcta según lo que realmente se
necesita: aquí no hace falta el orden completo, solo el \texttt{k}-ésimo
valor, así que ordenar todo es desperdiciar trabajo. Vale la pena mencionar
la alternativa de \emph{quickselect} (una variante de \emph{quicksort} que
solo particiona el lado que le interesa), que logra $O(n)$ en promedio pero
$O(n^2)$ en el peor caso, contra el $O(n \log k)$ garantizado del heap. Ten
cuidado con la trampa de usar un max-heap de todo el arreglo: funciona, pero
gasta $O(n)$ de espacio y $O(n + k \log n)$ de tiempo, peor que la versión
con heap de tamaño fijo \texttt{k}.
